\documentclass[10pt,a4paper]{amsart}
\usepackage[margin=19mm]{geometry}
\usepackage{amsmath,amssymb,mathtools}
\usepackage[scr=rsfs,cal=euler]{mathalfa}
\usepackage{xfrac}
\usepackage[unicode,hidelinks,bookmarks=false]{hyperref}
\newcommand{\ie}{i.e.\ }
\newcommand{\cf}{cf.\ }
\newcommand{\resp}{respectively\ }
\newcommand{\Subsection}{\subsection}
\providecommand{\nobreakdash}{\nobreak\hbox{-}}
\setlength{\emergencystretch}{2em}
\multlinegap0pt
\usepackage{amssymb}
\usepackage{bm}
\newtheorem{thm}{Theorem}
\newtheorem{cor}{Corollary}
\title{Geometric Shafarevich boundedness conjecture for families of polarized varieties}
\author{Junchao Shentu}
\date{Source: arXiv:2605.09847v1 (CC BY 4.0)}
\begin{document}
\maketitle

\noindent\emph{Source and licence.} Geometric Shafarevich boundedness conjecture for families of polarized varieties. Version arXiv:2605.09847v1 (CC BY 4.0), distributed by arXiv under the Creative Commons Attribution 4.0 International licence, http://creativecommons.org/licenses/by/4.0/, as recorded in the arXiv licence metadata for this exact version and shown on the abstract page https://arxiv.org/abs/2605.09847v1. Full licence notice: \texttt{licenses/arxiv-2605.09847-CC-BY-4.0.txt}.\par\smallskip

\noindent\textit{Editorial scope.} Counted unit: the article's cor cor\_Arakelov\_dim1 (source0.tex lines 844--849). The article's complete original proof of it is delivered with it (source0.tex lines 850--858, one \texttt{proof} environment of 9 lines). No mathematics is written, repaired or re-derived: the statement and the proof are the article's own lines, in the article's own notation, and no retained line is substituted or re-flowed.  Retained uncounted as the source-local context the proof needs: lines 684--699.  Nothing was omitted from any retained span, so the package carries no omission record.  No retained line contains a citation, so the wrapper carries no bibliography and no bibliography entry.  No retained line contains a figure environment, an image-inclusion command or drawing code, so no image is delivered and the package declares no asset dependency.  Editorial formatting only, in addition: an AMS \texttt{amsart} preamble that restates the article's own declarations and macros used by the retained lines, namely the environments \texttt{thm}, \texttt{cor} on separate counters, together with the standard packages those lines need.  The retained environments print in this document as Theorem 1, Corollary 1 in order of appearance, whereas the article prints the same items with the numbering of its own full document.  The complete native source of the article as distributed in the arXiv e-print for this exact version is bundled unchanged as \texttt{sources/200266/source0.tex}.
	\begin{thm}\label{thm_Arakelov_family}
		Let $f^o \colon (X^o, B^o) \to S^o$ be a birationally admissible family of $(d, \Phi_c, \Gamma, \sigma)$-stable minimal models over a smooth quasi-projective $n$-fold $S^o$, inducing a morphism $\xi^o \colon S^o \to M_{\mathrm{slc}}(d, \Phi_c, \Gamma, \sigma)$. Let $S$ be a smooth projective compactification of $S^o$ such that $D := S \setminus S^o$ is a reduced simple normal crossing divisor and $\xi^o$ extends to a morphism $\xi \colon S \to M_{\mathrm{slc}}(d, \Phi_c, \Gamma, \sigma)$. Let $(a,r,j)$ be a $(d, \Phi_c, \Gamma, \sigma)$-polarization datum. Define  
		\[
		\mathrm{ind}_\xi(\lambda_{a,r}) := \min\big\{k \in \mathbb{Z}_{>0} \,\big|\, \xi^*\lambda_{a,r}^{\otimes k} \text{ is a line bundle on } S \big\}.
		\]  
		If $K_S + D$ is pseudo-effective, then for every movable curve class $\alpha \in N_1(S)$ and every positive integer $k$ divisible by $\mathrm{ind}_\xi(\lambda_{a,r})$, the inequality  
		\begin{align}\label{align_Ara_1}
			c_1\big(\xi^*\lambda_{a,r}\big) \cdot \alpha 
			\leq \frac{1}{k(k l r d + 1)} \sum_{p=1}^{k l r d} p n^{p-1} (K_S + D) \cdot \alpha + \frac{2}{k}\, D \cdot \alpha,
		\end{align}  
		holds, where $l = \operatorname{rank}(\Lambda_{a,r})$.
		If moreover $K_S+D$ is ample, then  
		\begin{align}\label{align_Ara_2}
			c_1(\xi^*\lambda_{a,r})\cdot(K_S+D)^{n-1} \leq \frac{lrd}{n}(K_S+D)^{n}.
		\end{align}
	\end{thm}

\begin{cor}\label{cor_Arakelov_dim1}
	Let the notations and assumptions be as in Theorem \ref{thm_Arakelov_family}. Suppose further that $\dim S = 1$. Then  
	\begin{equation*}
		\deg\, \xi^\ast(\lambda_{a,r}) \leq \frac{r n \cdot \operatorname{rank}(\Lambda_{a,r})}{2} \cdot \deg(K_S + D).
	\end{equation*}
\end{cor}

\begin{proof}
	By Theorem \ref{thm_Arakelov_family}, for every positive integer $k$ divisible by $\operatorname{ind}_S(\lambda_{a,r})$, we have  
	\begin{equation*}
		\deg\, c_1\big(\xi^\ast(\lambda_{a,r})\big) 
		\leq \frac{1}{k(k l r n + 1)} \sum_{p=1}^{k l r n} p \cdot \deg(K_S + D) + \frac{2}{k} \deg D
		= \frac{l r n}{2} \deg(K_S + D) + \frac{2}{k} \deg D,
	\end{equation*}
	where $l := \operatorname{rank}(\Lambda_{a,r})$. Letting $k \to +\infty$, the term $\frac{2}{k}\deg D$ vanishes, yielding the claimed inequality.
\end{proof}

\end{document}
